% Graduate proof template. Copy the whole folder:
%   cp -r shared/templates/graduate-proof content/graduate-proofs/<slug>
% Label prefixes:  def:  lem:  thm:  cor:  prop:  eq:
% Refer with \cref{lem:first} ("Lemma 4.1") or \eqref{eq:name} ("(2.1)").
% Environments: theorem lemma proposition corollary definition example remark exercise proof
% Delete any section you don't need.
\documentclass[11pt,a4paper]{article}
\usepackage{dsmorgan}

\title{Theorem name}
\subtitle{One-line description}
\author{Drew M. Schaffner}
\date{Draft, \today}

\begin{document}
\maketitle

\section{Statement}

\begin{theorem}[Theorem name]\label{thm:main}
  State the result precisely, with every hypothesis.
\end{theorem}

\section{Setting}

What we assume, and the definitions the proof depends on.

\begin{definition}[Term]\label{def:first}
  Definition.
\end{definition}

\section{Idea of the proof}

Two or three sentences a reader can hold in their head before the details.

\section{Proof}

\begin{lemma}[Supporting fact]\label{lem:first}
  Statement.
\end{lemma}

\begin{proof}
  Proof of \cref{lem:first}.
\end{proof}

\begin{proof}[Proof of \cref{thm:main}]
  Citing \cref{def:first} and \cref{lem:first}.
\end{proof}

\section{Remarks}

\begin{remark}
  Consequences, generalisations, or where the hypotheses are really needed.
\end{remark}

\begin{thebibliography}{9}
  \bibitem{key} A.~Author, \emph{Title}, Publisher, Year.
\end{thebibliography}

\end{document}
